\section{Introducción}

\subsection{Problemas identificados}

Ante la crisis estructural del sistema de salud colombiano, se identificaron tres problemas principales que podrían beneficiarse de soluciones tecnológicas basadas en inteligencia artificial:

\subsubsection{Problema 1: Desabastecimiento de medicamentos y falta de alternativas}

Los pacientes enfrentan largas esperas y frecuentemente no encuentran los medicamentos prescritos en las farmacias. La falta de información sobre alternativas farmacéuticas con componentes activos similares genera frustración, abandono de tratamientos y deterioro de la salud, especialmente en personas con enfermedades crónicas.

\subsubsection{Problema 2: Ineficiencia en la gestión de inventarios farmacéuticos}

Las farmacias y gestores farmacéuticos carecen de sistemas predictivos que permitan anticipar desabastecimientos y optimizar la distribución de medicamentos. La falta de visibilidad en tiempo real sobre la disponibilidad de fármacos en diferentes puntos de la cadena de suministro genera pérdidas económicas y afecta la atención al paciente.

\subsubsection{Problema 3: Asimetría de información entre pacientes y el sistema de salud}

Los pacientes tienen acceso limitado a información confiable sobre sus medicamentos, alternativas disponibles, precios comparativos y disponibilidad en diferentes farmacias. Esta brecha informativa los deja en desventaja frente a un sistema de salud complejo y en crisis.

\subsection{Justificación como proyecto de Principios y Tecnologías de Inteligencia Artificial}

De los tres problemas identificados, se seleccionó el \textbf{Problema 1: Desabastecimiento de medicamentos y falta de alternativas} porque es ideal para un proyecto de PTIA por las siguientes razones:

\begin{itemize}
	\item \textbf{Requiere visión por computadora:} La identificación de medicamentos a partir de fotografías es un problema clásico de clasificación de imágenes que se beneficia de Redes Neuronales Convolucionales (CNN).
	
	\item \textbf{Involucra sistemas de recomendación:} Encontrar medicamentos similares basándose en componentes activos y características extraídas requiere algoritmos de aprendizaje automático como K-Nearest Neighbors (KNN).
	
	\item \textbf{Combina múltiples técnicas de IA:} El proyecto integra deep learning (CNN para clasificación) con machine learning tradicional (KNN para recomendación), permitiendo explorar diferentes enfoques y compararlos.
	
	\item \textbf{Permite aprendizaje supervisado:} Se puede entrenar el modelo con un dataset etiquetado de medicamentos, aplicando técnicas de data augmentation, transfer learning y fine-tuning.
	
	\item \textbf{Tiene métricas evaluables:} Es posible medir el desempeño del sistema con métricas estándar como accuracy, precision, recall y F1-score para clasificación, y métricas de similitud para recomendación.
	
	\item \textbf{Escalabilidad y mejora continua:} El modelo puede mejorarse continuamente con más datos, explorando arquitecturas más avanzadas como ResNet, EfficientNet o Vision Transformers.
\end{itemize}

\subsection{Contexto de la crisis del sistema de salud colombiano}

En los últimos años se ha venido presentando una problemática país con respecto al sector de salud, específicamente hacia el sector de las EPS (Entidad Promotora de Salud), en donde hace dos años el presidente vigente Gustavo Petro denunció reportes fraudulentos sobre la distribución de sus recursos para la atención en salud y fue el año pasado que, según informes presentados por el Ministerio de Salud el 15 de Julio del 2025, comenta que esa deuda (en el cual reúne al menos 29 EPS) es de aproximadamente 35 billones de dólares \cite{minsalud2025}. Esta deuda genera una crisis en el sistema de salud porque el gobierno está acorralando y asfixiando a una posible liquidez de las IPS (Instituciones Prestadoras de Servicios de Salud), provocando el cierre de servicios, falta de pagos a médicos, desabastecimiento de medicamentos y barreras de acceso para pacientes. La problemática financiera desencadena una serie de conflictos, no solamente en la sostenibilidad de hospitales y clínicas, como demoras y falta de medicamentos.

Bajo esta naciente preocupación, donde aumenta de sobremanera los tiempos de espera para su distribución, afectando la continuidad de los tratamientos, en el último año se han venido haciendo cada vez más evidentes. Acorde a una investigación por parte de la Corte Institucional, nos demuestra la realidad, la cual no es nada favorable:

\begin{quote}
	``La Corte Constitucional constató que persisten graves dificultades en la entrega de fármacos por deudas acumuladas entre EPS y gestores farmacéuticos, incremento de tutelas de pacientes, retrasos en trámites del Invima y ausencia de información confiable para alertar los desabastecimientos.'' \cite{colprensa2025}
\end{quote}

Se han ido proponiendo e implementando estrategias desde que inició el año para mitigar los daños, aparte de los ya presentes, entre ellas que el estado se haga cargo de los recursos de las EPS y que estas ultimas se encarguen solamente de prestar un servicio únicamente de gestores de vida y salud, además de un fuerte aumento en el presupuesto de los recursos destinados al sector salud.

Sin embargo, toda esta problemática sobre el asunto de los medicamentos no es algo que podrá encontrar una solución de manera inmediata porque a pesar de que estalló con la deuda de las EPS, es algo que venía de años atrás. Aún así, la sociedad colombiana no se puede quedar esperando por una respuesta positiva por parte de las entidades de salud tanto tiempo, porque la brecha que irá dejando a su paso será inmensa.

No es un tema sencillo de resolver, pero en unión con las innovaciones y revoluciones utilizando la tecnología como medio, se puede encontrar aquel punto óptimo, así mientras se termina de dar el camino más adecuado para el sector salud en Colombia, la sociedad colombiana no termine siendo los más afectados ante todo este conflicto.

\pagebreak